\documentclass[11pt]{article}

\usepackage[margin=1in]{geometry}
\usepackage{booktabs}
\usepackage{amsmath}
\usepackage{amssymb}
\usepackage{hyperref}
\usepackage{cite}
\usepackage{authblk}

\title{Deciphering Molecular Heterogeneity and Dynamics of Human Hippocampal Neural Stem Cells from Infancy to Aging and Stroke Injury}

\author[1]{Anonymous Author}
\affil[1]{Anonymous Institution}

\date{}

\begin{document}

\maketitle

\begin{abstract}
Whether adult hippocampal neurogenesis persists in humans, and whether neural stem cells (NSCs) can be reactivated after brain injury, remains controversial. Here we use droplet-based single-nucleus RNA sequencing (snRNA-seq) of hippocampi from $10$ human donors spanning neonatal to aged life and including a stroke case, together with hippocampus from a 3-month-old cynomolgus monkey, to map the molecular heterogeneity and dynamics of the human hippocampal neurogenic lineage. Samples were grouped by age and health into neonatal (Day 4 after birth; $n = 1$), adult (31- and 32-year-old; $n = 2$), aging (50- to 68-year-old; $n = 6$), and stroke-injured (48-year-old; $n = 1$). After excluding cells with $>20\%$ mitochondrial-gene transcripts, we recover a median of $3001$ genes per nucleus. We identify five molecularly distinct neurogenic populations --- qNSC1, qNSC2, primed NSC (pNSC), active NSC (aNSC) and neuroblast/immature granule cell (NB/im-GC) --- each with specific marker signatures; pNSC and aNSC are abundant in neonatal hippocampus and are re-activated by ischaemic injury. In the stroke case, pNSC and aNSC populations reach $17.3\%$ ($322/1861$) and $11.7\%$ ($218/1861$) of the neurogenic lineage, comparable to neonatal ratios and significantly above adult and aged levels. Reciprocally, qNSC populations in the injury group fall to $23.4\%$ (qNSC1, $413/1258$) and $11.1\%$ (qNSC2, $214/1258$), compared with adult ($32.8\%$ / $17.0\%$) and aged ($57.8\%$ / $22.6\%$). Monocle pseudotime trajectories, branched-expression analysis modelling (BEAM), and single-cell hierarchical Poisson factorization (scHPF) support a continuous trajectory from qNSC $\to$ pNSC $\to$ aNSC $\to$ NB/im-GC, and the re-integrative analysis with Zhou et al.'s $14$-donor aged snRNA-seq dataset recovers no evident pNSC/aNSC/NB populations in uninjured aged donors. Immunostaining with HOPX, NES, KI67, and PSA-NCAM confirms the transcriptomic classification in situ.
\end{abstract}

\section{Introduction}

Whether adult hippocampal neurogenesis persists in humans, and to what extent, has been debated for a quarter of a century \cite{eriksson1998,spalding2013,sorrells2018,boldrini2018,moreno2019,kempermann2018human}. Two influential studies \cite{sorrells2018,moreno2019} report that neurogenic activity drops to undetectable levels in healthy adult hippocampus, whereas others document sustained neurogenesis across the lifespan \cite{boldrini2018,spalding2013}. Methodologically, immunostaining-based detection of rare neurogenic cells is sensitive to tissue-processing artefacts, while conventional RNA-seq cannot discriminate the $\le 1\%$ rare cell populations involved. Recent advances in single-nucleus RNA sequencing (snRNA-seq) \cite{habib2017nuclei,lacar2016nuclear} provide a third route, allowing thousands of hippocampal nuclei to be profiled simultaneously from frozen post-mortem tissue \cite{zhou2022snrnaseq}.

A second open question is whether the adult human neurogenic lineage can be re-activated by brain injury. Rodent studies have shown that stroke recruits a quiescent pool into active neurogenesis \cite{arvidsson2002stroke,jin2006stroke,duan2014stroke}, and that reactive astrogliosis and neurogenesis share overlapping molecular programs \cite{liddelow2017neurotoxic,magnusson2014regen,kimelberg2010reactive}. Whether a similar programme operates in humans has not been directly assayed. Here we profile hippocampi from $10$ donors spanning neonatal to aged life, including a stroke case, using snRNA-seq, and integrate the result with a hippocampal snRNA-seq atlas from $14$ aged donors \cite{zhou2022snrnaseq}. We resolve five neurogenic populations along a continuous trajectory and show that ischaemic injury dramatically re-activates the primed and active NSC compartments.

\section{Methods}

\paragraph{Donors.} We recruited $10$ donors from neonatal day 4 to $68$ years old, grouped as follows: neonatal (postnatal $4$ days, $n = 1$), adult (31y, 32y, $n = 2$), aging (50y, 56y, 60y, 64y-1, 64y-2, 68y, $n = 6$), and stroke injury (48y, $n = 1$); $1$ female and $9$ males. Given the differences between the rostral and caudal hippocampus \cite{wu2014hippocampus}, we used the anterior (AN) and mid (MI) hippocampus containing typical dentate gyrus (DG) structures for snRNA-seq and immunostaining. A female cynomolgus monkey aged $3$ months (body weight $2.3$\,kg) was used for cross-species validation. Samples were grouped into four sets: neonatal (Day $4$ after birth, D4; $n=1$), adult (31--32-year-old, $n=2$), aging (50--68-year-old, $n=6$), and stroke-induced injury (48-year-old, $n=1$) (Fig.~1A).

\paragraph{Nucleus isolation.} Tissue samples were frozen in OCT (Tissue-Tek) on dry ice and sectioned at $10\,\mu$m on a cryostat (Leica CM1950). Chilled lysis buffer ($10$\,mM Tris-HCl, $10$\,mM NaCl, $3$\,mM MgCl$_2$, and $0.1\%$ Nonidet P40 Substitute; Sigma Aldrich, PN-74385) was added to the tissue and the tube was incubated on ice for $15$\,min with gentle shaking. A $30\,\mu$m MACS SmartStrainer (Miltenyi Biotec, PN-130-098-458) was used to remove cell debris and large clumps. Nuclei were centrifuged at $500{\times}g$ for $5$\,min at $4\,^\circ$C, washed in $1\times$ PBS with $1.0\%$ BSA and $0.2$\,U/$\mu$L RNase inhibitor (Sigma Aldrich, PN-3335399001), and gently pipetted $8$--$10$ times.

\paragraph{snRNA-seq and processing.} Libraries were prepared on the 10x Genomics Chromium platform \cite{zheng201710x}. After removal of cell debris, cell aggregates, and cells with $>20\%$ mitochondrial-gene transcripts, a median of $3001$ genes per nucleus was retained. Seurat \cite{stuart2019seurat3,butler2018seurat,hao2021sc} was used for normalization, dimensionality reduction, clustering, and differential expression. Marker genes were identified by Wilcoxon rank-sum test with average expression difference $>0.5$ (natural log) and $p < 0.05$.

\paragraph{Pseudotime and scHPF.} Monocle 2 \cite{trapnell2014monocle,qiu2017monocle2} pseudotime was computed using dispersed genes ($\mathrm{estimateDispersions}$), the DDRTree dimensionality reduction (\texttt{reduceDimension(mycds, max\_components = 2, method = 'DDR Tree')}), and $\mathrm{plot\_cell\_trajectory}$. Branched-expression analysis modelling (BEAM) was used to identify branch-specific genes; differentiation-related DEGs were called at $q < 1 \times 10^{-4}$ and partitioned into four gene clusters. scHPF \cite{levitin2019schpf} was run in parallel with Seurat \texttt{FindAllMarkers} to identify new marker-gene signatures for the neurogenic lineage, and gene-ontology analysis was performed using \texttt{clusterProfiler} \cite{yu2012clusterprofiler}.

\paragraph{Immunohistochemistry.} Human and macaque hippocampal tissue was fixed in $4\%$ paraformaldehyde (PFA) for up to $24$\,h and cryoprotected in $30\%$ sucrose at $4\,^\circ$C until completely sunk. Tissue slides (anterior hippocampus containing typical DG) were incubated in blocking/permeation solution with $2\%$ Triton X-100 (Sigma) for $2$\,h, then treated with a VECTOR TrueView autofluorescence-quenching kit (Vectorlabs, PN-SP-8500-15), washed $3 \times 15$\,min in PBS (pH $7.6$), and incubated in $3\%$ BSA for $1$\,h. Sections were incubated overnight with primary antibodies at $4\,^\circ$C, washed $3 \times 15$\,min in PBS, and incubated for $2$\,h at room temperature with Alexa Fluor 488 AffiniPure donkey anti-rabbit IgG (H+L) ($1{:}500$; Jackson ImmunoResearch, PN-712-545-152), Alexa Fluor 647 AffiniPure donkey anti-rabbit IgG (H+L) ($1{:}500$; Jackson ImmunoResearch, PN-715-605-150), or Alexa Fluor 555 ($1{:}500$; ThermoFisher, PN-A-31572). DAPI (Sigma, PN-32670-5mg-F) was used for counterstaining. TUNEL was performed by permeabilizing in $0.02\%$ Triton X-100 overnight, washing, adding $20\,\mu$L of TUNEL reaction mixture (labelling buffer, TdT, BIO-d-UTP) for $120$\,min at $37\,^\circ$C in a humidified chamber, washing with PBS for $2$\,min, and blocking with $50\,\mu$L blocking reagent at room temperature for $30$\,min.

\section{Results}

\subsection{A snRNA-seq atlas of human hippocampus across life and stroke}
We profiled $10$ human hippocampi (neonatal $n = 1$; adult $n = 2$; aging $n = 6$; stroke $n = 1$) and one cynomolgus monkey hippocampus. After quality filtering, a median of $3001$ genes per nucleus was retained (Fig.~1-figure supplement 1A--B). Seurat clustering resolved the full complement of hippocampal cell types, and focused analysis of neurogenic populations identified qNSC1, qNSC2, pNSC, aNSC, and NB/immature granule cell (NB/im-GC) clusters, characterized by markers including HOPX and LPAR1 (radial glia-like), S100B and GFAP (astrocyte), ST18, STMN1, PROX1, and SIRT2 (neuronal development).

\subsection{Continuous trajectory qNSC $\to$ pNSC $\to$ aNSC $\to$ NB/im-GC}
Monocle pseudotime placed the five neurogenic populations along a continuous trajectory, with HOPX, VIM, CHI3L1, and TNC expressed early (NSC genes) and NRGN, STMN2, and SV2B expressed late (neuronal genes). Cell-cycle scoring showed that pNSC and aNSC are enriched for S-/G2M-phase signatures, whereas qNSC1 and qNSC2 are predominantly G1-quiescent. scHPF factor analysis and Seurat \texttt{FindAllMarkers} converged on overlapping but distinct signatures for the pNSC and aNSC populations, and GO-term enrichment of the top 1000 pNSC and aNSC genes highlighted neural-development-related terms at $p < 0.05$. At 3 months in cynomolgus monkey hippocampus, the neuroblast marker DCX was also expressed in a subset of interneurons; we therefore filtered NB/im-GC genes against the interneuron population by intersecting the top 100 scHPF factors with $\mathrm{FindAllMarkers}$ at adjusted $p < 0.01$.

\subsection{Stroke injury re-activates the pNSC/aNSC compartment}
Quantification of qNSC1, qNSC2, pNSC, aNSC, reactive astrocytes, and NB across neonatal (N), adult (Ad), aging (Ag), and stroke-injured (I) hippocampi revealed that pNSC and aNSC populations in the injury group reached $17.3\%$ ($322/1861$) and $11.7\%$ ($218/1861$) of the neurogenic lineage, respectively --- comparable to neonatal ratios and significantly higher than in adult and aging groups. Correspondingly, qNSC1 and qNSC2 ratios in the injury group dropped to $23.4\%$ ($413/1258$) and $11.1\%$ ($214/1258$), compared with adult ($32.8\%$; $17.0\%$) and aged ($57.8\%$, $310/536$; $22.6\%$, $121/536$) groups (Table~\ref{tab:ratios}). Up-regulated genes in the injured neurogenic lineage were enriched for GO terms related to proliferation and regeneration at $p < 0.05$.

\begin{table}[h]
\centering
\caption{Percentages of qNSC1, qNSC2, pNSC, and aNSC within the neurogenic lineage across groups.}
\label{tab:ratios}
\begin{tabular}{lcccc}
\toprule
Group & qNSC1 & qNSC2 & pNSC & aNSC \\
\midrule
Adult & $32.8\%$ ($413/1258$) & $17.0\%$ ($214/1258$) & -- & -- \\
Aging & $57.8\%$ ($310/536$) & $22.6\%$ ($121/536$) & -- & -- \\
Stroke & $23.4\%$ ($413/1258$) & $11.1\%$ ($214/1258$) & $17.3\%$ ($322/1861$) & $11.7\%$ ($218/1861$) \\
\bottomrule
\end{tabular}
\end{table}

\subsection{Pseudotime in stroke-injured hippocampus}
Pseudotime reconstruction of the stroke-injured neurogenic lineage again placed pNSCs and aNSCs between qNSC populations and NB/im-GC, with reactive astrocytes positioned adjacent to the primed state, consistent with rodent models showing that reactive astrocytes and primed NSCs share molecular programs \cite{magnusson2014regen,liddelow2017neurotoxic,kimelberg2010reactive}. Neonatal-day-4 Monocle trajectories showed that pNSCs generate two branches (GC1 and GC2), with BEAM identifying four gene clusters at $q < 1 \times 10^{-4}$.

\subsection{Re-integration with an aged snRNA-seq atlas}
We integrated our snRNA-seq with Zhou's snRNA-seq dataset of $14$ aged donors (from $60$--$92$ years old). In the integrated object we did not detect evident pNSC, aNSC, or NB populations in the $14$ Zhou aged donors (Fig.~6-figure supplement 11B), consistent with the decline of the primed/active compartment with age and its re-emergence after injury. Immunostaining for HOPX/NES, KI67/NES, and PSA-NCAM (scale bars $200\,\mu$m, $100\,\mu$m, and $10\,\mu$m) corroborated the transcriptomic classifications in situ.

\section{Discussion}

By profiling $10$ human hippocampi spanning infancy to aging --- with one stroke case --- at the single-nucleus level, we resolve five molecularly distinct neurogenic populations and map them onto a continuous trajectory from qNSC $\to$ pNSC $\to$ aNSC $\to$ NB/im-GC. Quantitatively, the primed and active NSC pools are abundant in neonatal hippocampus, sharply contract in adult and aged donors, and rebound dramatically after stroke (pNSC $= 17.3\%$ of lineage; aNSC $= 11.7\%$), comparable to neonatal ratios and significantly above adult and aging levels. Reciprocally, the quiescent pool contracts (qNSC1: $23.4\%$; qNSC2: $11.1\%$ in injury vs.\ adult $32.8\%$ / $17.0\%$; aged $57.8\%$ / $22.6\%$). Re-integration with Zhou et al.'s snRNA-seq of $14$ aged donors does not recover detectable pNSC/aNSC/NB populations, indicating that this injury-induced reactivation is not an artefact of cohort differences.

This work addresses both of the two long-standing controversies in the field. First, it provides molecular evidence for ongoing neurogenic potential in the adult human hippocampus that is masked in the quiescent state and revealed only when rare transitional (pNSC, aNSC, NB/im-GC) populations are sampled deeply. Second, it demonstrates that ischaemic injury can re-activate this latent compartment, recapitulating the neonatal transcriptomic program. Limitations include the small sample size in the neonatal ($n = 1$) and stroke ($n = 1$) groups, the limited information about patients' health status and lifestyle parameters, and the cross-sectional nature of the design. Nevertheless, the convergent evidence from snRNA-seq, scHPF factorization, Monocle pseudotime with BEAM, in situ immunostaining, and cross-species (macaque) comparison supports a coherent model in which a rare but persistent neurogenic compartment is maintained in human hippocampus throughout life and is recruitable by brain injury. These findings motivate therapeutic strategies that harness the latent neurogenic potential of the adult human hippocampus for neuronal replacement after stroke and other acute injuries.

\bibliographystyle{unsrt}
\bibliography{references}

\end{document}
